\documentclass[12pt,a4paper]{article}

\usepackage{amsmath, amssymb, amsthm}
\usepackage{mathtools}
\usepackage{geometry}
\usepackage{hyperref}
\usepackage{booktabs}
\usepackage{array}
\usepackage{xcolor}

\geometry{margin=2.5cm}

\newtheorem{remark}{Remark}
\newtheorem{proposition}{Proposition}

\newcommand{\R}{\mathbb{R}}
\newcommand{\cL}{\mathcal{L}}
\newcommand{\cG}{\mathcal{G}}
\newcommand{\cD}{\mathcal{D}}
\newcommand{\cH}{\mathcal{H}}
\newcommand{\cM}{\mathcal{M}}
\newcommand{\tr}{\operatorname{tr}}
\newcommand{\SVD}{\operatorname{SVD}}
\newcommand{\diag}{\operatorname{diag}}
\newcommand{\inner}[2]{\langle #1,\, #2 \rangle}

\title{Riemannian Steepest Descent on the Fixed-Rank Matrix Manifold\\
\large Derivation and Algorithm}
\author{}
\date{}

\begin{document}

\maketitle

% -----------------------------------------------------------------------
\section{Problem Statement}

Let $n, r \in \mathbb{N}$ with $r \ll n$. We consider the problem of finding a
rank-$r$ matrix $X \in \R^{n \times n}$ that satisfies the generalized
Sylvester-like equation
\begin{equation}\label{eq:original}
    \cL[X] := A_1 X + X A_2 + A_3 X A_4 \;=\; Y,
\end{equation}
where $A_i \in \R^{n \times n}$ are given matrices and $\mathcal{L}$ is self-adjoint. We enforce the
rank constraint via the parametrization
\begin{equation}\label{eq:factorization}
    X = U S V^\top, \qquad U, V \in \R^{n \times r},\quad
    U^\top U = V^\top V = I_r,\quad S \in \R^{r \times r},
\end{equation}
with $U$ and $V$ having orthonormal columns. The constraint set is the \emph{fixed-rank manifold}
\begin{equation}
    \cM_r \;=\; \bigl\{X \in \R^{n \times n} : \operatorname{rank}(X) = r\bigr\}.
\end{equation}

Equation~\eqref{eq:original} is rewritten as the condition that the
\emph{defect operator}
\begin{equation}\label{eq:operator}
    \cG[X] \;=\; A_1X + XA_2 + A_3 X A_4 - Y
\end{equation}
vanishes. Because we restrict to $\cM_r$, we cannot enforce $\cG[X] = 0$
in the ambient space; instead we impose the \emph{projected} condition
\begin{equation}\label{eq:projected}
    P(X)\,\cG[X] \;=\; 0,
\end{equation}
where $P(X)$ is the orthogonal projector onto the tangent space of $\cM_r$
at $X$, given as
\begin{align*}
    P(X)\,Z \;=\; U U^\top Z (I - V V^\top) + Z V V^\top\,.
\end{align*}
Equation \eqref{eq:projected} is the first-order optimality condition of the quadratic cost
\begin{equation}\label{eq:cost}
    \Phi(X) \;=\; \tfrac{1}{2}\,\inner{\cL[X] - Y}{X}
\end{equation}
under the constraint $X\in\mathcal{M}_r$ where $\|\cdot\|_F$ denotes the Frobenius norm and
$\inner{A}{B} = \tr(A^\top B)$ is the associated inner product. We assume that the Hessian of $\Phi$ is \emph{symmetric positive definite
(SPD)} as a linear operator on $\R^{n \times n}$.

% -----------------------------------------------------------------------
\section{DLRT with Exact Line Search}

We minimize $\Phi$ restricted to $\cM_r$ using Riemannian steepest
descent.  The Riemannian gradient at $X_k \in \cM_r$ is the tangent vector
obtained by projecting the Euclidean gradient onto $T_{X_k}\cM_r$:
\begin{equation}
    \operatorname{grad}_{\cM_r}\Phi(X_k) \;=\; P(X_k)\,\nabla\Phi(X_k)
    \;=\; P(X_k)\,\cG[X_k].
\end{equation}
Steepest descent moves in the direction of the negative Riemannian gradient.

\subsection{Step 1 — Residual (Riemannian Gradient)}

Define the \emph{residual} as the negative Riemannian gradient:
\begin{equation}\label{eq:residual}
    \boxed{R_k \;=\; -P(X_k)\,\cG[X_k].}
\end{equation}
By construction $R_k \in T_{X_k}\cM_r$.  The Frobenius norm $\|R_k\|_F$
measures how far $X_k$ is from satisfying the projected equation and
serves directly as the convergence criterion.

\subsection{Step 2 — Search Direction}

The steepest-descent direction in the Riemannian sense is
\begin{equation}\label{eq:direction}
    \boxed{D_k \;=\; R_k.}
\end{equation}
Since $R_k \in T_{X_k}\cM_r$, the direction is already a valid tangent
vector and requires no further projection.

\subsection{Step 3 — Optimal Step Size (Exact Line Search)}

We perform an exact line search along the tangent direction $D_k$, i.e.\ we
minimize the restriction of $\Phi$ to the curve $\alpha \mapsto X_k + \alpha D_k$
(before retraction).  Since $\Phi$ is quadratic, this one-dimensional
problem has the closed-form minimizer
\begin{equation}\label{eq:stepsize-derivation}
    \frac{d}{d\alpha}\Phi(X_k + \alpha D_k)\Big|_{\alpha=\alpha_k}
    = \inner{\nabla\Phi(X_k + \alpha_k D_k)}{D_k}
    = \inner{\cG[X_k] + \alpha_k \cL[D_k]}{D_k}
    = 0.
\end{equation}
Solving for $\alpha_k$ while using $D_k = R_k = -P(X_k)\cG[X_k]$ and $\langle \cG[X_k], D_k \rangle = \langle \cG[X_k], P(X_k)D_k \rangle=\langle P(X_k)\cG[X_k], D_k \rangle$ gives
\begin{align}
    \alpha_k = -\frac{\inner{\cG[X_k]}{D_k}}{\inner{\cL[D_k]}{D_k}} = -\frac{\inner{P(X_k)\cG[X_k]}{D_k}}{\inner{\cL[D_k]}{D_k}} = \frac{\inner{R_k}{R_k}}{\inner{\cL[D_k]}{D_k}}\,.
\end{align}

\subsection{Step 4 — Retraction onto $\cM_r$}

The updated point $X_k + \alpha_k D_k$ generally has rank up to $2r$ and
therefore lies \emph{off} the manifold $\cM_r$.  We map it back via the
\emph{rank-$r$ truncated SVD retraction}. This can be done with a DLRT update.


% -----------------------------------------------------------------------
\section{Complete Algorithm}

\begin{center}
\renewcommand{\arraystretch}{1.5}
\begin{tabular}{ll}
\toprule
\textbf{Input}  & Initial guess $X_0 = U_0 S_0 V_0^\top \in \cM_r$,
                  tolerance $\varepsilon > 0$. \\
\textbf{Output} & Approximate solution $X_k \in \cM_r$ with
                  $\|R_k\|_F \le \varepsilon$. \\
\bottomrule
\end{tabular}
\end{center}

\medskip
\noindent For $k = 0, 1, 2, \ldots$:

\begin{enumerate}
    \item \textbf{Define the gradient:}
    \[
        \cG[X_k] = A_1X_k + X_kA_2
                   + A_3 X_k A_4 - Y
    \]
    and \textbf{compute the projected residual / Riemannian gradient:}
    \[
        R_k = -\Bigl[U_k U_k^\top \cG[X_k](I{-}V_k V_k^\top)
              + \cG[X_k]\,V_k V_k^\top\Bigr].
    \]
    The residual lies in the tangent space, so it can be stored as the three factors $ U_k^\top\mathcal{L}[X_k]$, $\mathcal{L}[X_k]V_k $, and $U_k^\top\mathcal{L}[X_k]V_k$, which we can compute efficiently.

    \item \textbf{Check convergence:} if $\|R_k\|_F^2 = \| U_k^\top\mathcal{G}[X_k]\|_F^2 + \| \mathcal{G}[X_k]V_k\|_F^2 - \| U_k^\top\mathcal{G}[X_k]V_k \|_F^2 \le \varepsilon$, stop.

    \item \textbf{Set search direction:} $D_k = R_k$.

    \item \textbf{Define the Hessian on the direction:}
    \[
        \cH[D_k] := \cL[D_k] = A_1D_k + D_kA_2
                   + A_3 D_k A_4
    \]
    and \textbf{compute the projected Hessian:}
    \begin{align*}
        P(X_k)\cH[D_k] =\,& U_k U_k^\top \cH[D_k](I{-}V_k V_k^\top)
                   + \cH[D_k]\,V_k V_k^\top\\
                   =\,& U_k U_k^\top \cH[D_k]
                   + (I{-}U_k U_k^\top)\cH[D_k]\,V_k V_k^\top
    \end{align*}
    The projected Hessian lies in the tangent space, so it can be stored efficiently through the three factors $U_k^\top\cH[D_k]$, $\cH[D_k]V_k$, and $U_k^\top\cH[D_k]V_k $.

    \item \textbf{Compute the optimal step size:}
    \begin{align*}
        \alpha_k =\,& \frac{\|R_k\|_F^2}{\langle P(X_k)D_k,\cH[D_k]\rangle}\\
        =\,& \frac{\| U_k^\top\mathcal{G}[X_k]\|_F^2 + \| \mathcal{G}[X_k]V_k\|_F^2 - \| U_k^\top\mathcal{G}[X_k]V_k \|_F^2}{\langle D_k,P(X_k)\cH[D_k]\rangle}
    \end{align*}
    which can be computed efficiently with
    \begin{align*}
        \langle D_k,P(X_k)\cH[D_k]\rangle =\,& \langle -\Bigl[U_k U_k^\top \cG[X_k](I{-}V_k V_k^\top)
              + \cG[X_k]\,V_k V_k^\top\Bigr],P(X_k)\cH[D_k]\rangle \\
              =\,& -\langle U_k U_k^\top \cG[X_k](I{-}V_k V_k^\top),P(X_k)\cH[D_k]\rangle - \langle \cG[X_k]\,V_k V_k^\top,P(X_k)\cH[D_k]\rangle \\
              =\,& -\langle U_k^\top \cG[X_k](I{-}V_k V_k^\top),U_k^\top \cH[D_k]\rangle - \langle \cG[X_k]\,V_k ,\cH[D_k]\,V_k \rangle \\
              =\,& -\langle U_k^\top \cG[X_k],U_k^\top \cH[D_k]\rangle + \langle U_k^\top \cG[X_k]V_k,U_k^\top \cH[D_k]V_k\rangle- \langle \cG[X_k]\,V_k ,\cH[D_k]\,V_k \rangle
    \end{align*}

    \item \textbf{Dynamical low-rank update:}
    \begin{alignat*}{2}
        K_{k+1} =\,& U_kS_k - \alpha_k \mathcal{G}[X_k]V_k\,, \quad &&\widehat U = \textrm{orth}(K_{k+1})\\
        L_{k+1} =\,& V_kS_k^{\top} - \alpha_k (U_k^{\top}\mathcal{G}[X_k])^{\top}\,, \quad &&\widehat V = \textrm{orth}(L_{k+1})\\
        \widehat S =\,& (\widehat U^{\top}U_k)S_k(V_k^{\top}\widehat V) - \alpha_k \widehat U^{\top}\mathcal{G}[X_k]\widehat V.
    \end{alignat*}
    Extract $U_{k+1}, S_{k+1}, V_{k+1}$ from the truncated SVD.
\end{enumerate}

\end{document}